\section{Problem Analysis and Restatement}
\label{sec:problem_analysis}

\subsection{Semantic Parsing and Problem Understanding}

The problem concerns optimizing delivery routes for humanitarian aid distribution under uncertain road network conditions. To develop an effective mathematical model, we first perform a systematic analysis of the problem structure, identifying key components, relationships, and constraints.

\textbf{System Boundaries.} The system encompasses:
\begin{itemize}
\item Spatial domain: A geographic region with $N$ distribution points (health centers, clinics, communities requiring supplies)
\item Temporal domain: A planning horizon (daily, weekly) with delivery time windows
\item Resource domain: Vehicles with capacity constraints, warehouses with inventory limits
\item Uncertainty domain: Road network conditions affecting travel times probabilistically
\end{itemize}

\textbf{System Participants.}
\begin{itemize}
\item \textit{Humanitarian Organization:} Decision-maker responsible for planning and executing deliveries
\item \textit{Warehouse Operators:} Manage inventory at central and local warehouses
\item \textit{Vehicle Drivers:} Execute delivery routes, must adhere to safety and operational constraints
\item \textit{Beneficiaries:} Populations at distribution points receiving life-saving supplies
\end{itemize}

\textbf{System Objects.}
\begin{itemize}
\item \textit{Warehouses:} Central depot (Task A) or multiple local depots (Task B) storing supplies
\item \textit{Vehicles:} Transportation resources with finite capacity and range
\item \textit{Roads:} Network links connecting locations, characterized by stochastic travel times
\item \textit{Supplies:} Vaccines, medicines, and medical equipment with volume, weight, and urgency requirements
\end{itemize}

\subsection{Variable Identification}

We classify variables into three categories: state variables, decision variables, and parameters.

\subsubsection{State Variables (Dynamic System Properties)}

\begin{itemize}
\item $S_{i}(t)$: Inventory level at warehouse $i$ at time $t$
\item $P_{ij}(t)$: Probability that road $(i,j)$ is passable at time $t$
\item $T_{ij}(t)$: Random travel time on road $(i,j)$ at time $t$
\item $V_k(t)$: Position of vehicle $k$ at time $t$
\item $D_i(t)$: Demand level at distribution point $i$ at time $t$
\end{itemize}

\subsubsection{Decision Variables (Control Actions)}

\begin{itemize}
\item $x_{ijk}$: Binary variable indicating if vehicle $k$ travels from node $i$ to node $j$
\item $y_{ik}$: Binary variable indicating if node $i$ is served by vehicle $k$
\item $t_i$: Arrival time at node $i$
\item $q_k$: Quantity of supplies loaded on vehicle $k$
\item $z_{ijk}$: Binary variable for road selection considering reliability
\end{itemize}

\subsubsection{Parameters (Fixed Problem Data)}

\begin{itemize}
\item $N$: Number of distribution points
\item $K$: Number of available vehicles
\item $C_k$: Capacity of vehicle $k$
\item $d_i$: Demand at distribution point $i$
\item $Q$: Warehouse inventory
\item $\mu_{ij}$: Mean travel time on road $(i,j)$
\item $\sigma_{ij}^2$: Variance of travel time on road $(i,j)$
\item $\alpha$: Minimum reliability threshold (e.g., $\alpha = 0.85$ for 85\% reliability)
\item $[e_i, l_i]$: Time window for delivery at node $i$
\end{itemize}

\subsection{Problem Restatement}

\subsubsection{Precise Formulation}

We now provide a precise mathematical restatement of the problem.

\textbf{Task A: Single-Warehouse Stochastic VRP with Reliability Constraints}

Given:
\begin{itemize}
\item A complete directed graph $G = (V, E)$ where $V = \{0, 1, 2, ..., N\}$ is the set of nodes (0 = central warehouse, $1, ..., N$ = distribution points)
\item $E = \{(i,j) : i, j \in V, i \neq j\}$ is the set of road links
\item Travel times $\tau_{ij}$ on each edge $(i,j)$ are random variables following probability distribution $F_{ij}(\cdot)$
\item Vehicle fleet $K$ with capacities $C_k$, $k = 1, ..., K$
\item Demands $d_i$ at each node $i \in V \setminus \{0\}$
\item Time windows $[e_i, l_i]$ for service at each node
\item Reliability requirement $\alpha \in (0,1]$
\end{itemize}

Determine:
\begin{itemize}
\item Routes $R_k = (0, i_1, i_2, ..., i_{n_k}, 0)$ for each vehicle $k \in K$
\item Service times $t_i$ for each node $i$
\end{itemize}

Such that:
\begin{enumerate}
\item All distribution points are served: $\bigcup_{k=1}^{K} \{i_1, ..., i_{n_k}\} = V \setminus \{0\}$
\item Vehicle capacity constraints: Total demand on each route $\leq C_k$
\item Time window constraints: $e_i \leq t_i \leq l_i$ for all $i$
\item Reliability constraints: $\Pr(\text{Route completed within time horizon}) \geq \alpha$
\item Route continuity and connectivity
\item Minimize expected total delivery time: $\mathbb{E}[\sum_{k=1}^{K} T(R_k)]$
\end{enumerate}

\textbf{Task B: Multi-Warehouse Stochastic VRP with Coordination}

Extension of Task A where:
\begin{itemize}
\item Multiple warehouses $W = \{w_1, w_2, ..., w_M\}$
\item Vehicles assigned to specific warehouses: $k \in K_w$ for warehouse $w$
\item Decision variables include:
\begin{itemize}
\item Partitioning of nodes among warehouses
\item Inter-warehouse inventory transfers
\item Coordinated departure times
\end{itemize}
\end{itemize}

\subsection{Modeling Assumptions}

To make the problem tractable while maintaining practical relevance, we make the following assumptions:

\textbf{Structural Assumptions:}
\begin{enumerate}
\item The road network is connected (a path exists between any two points)
\item Vehicles are homogeneous (same capacity) unless specified otherwise
\item Demands are deterministic and known in advance
\item Warehouse capacity is sufficient for total demand
\item Vehicles must return to their starting warehouse
\end{enumerate}

\textbf{Stochastic Assumptions:}
\begin{enumerate}
\item Travel times on different road segments are independent random variables
\item Travel time distributions are stationary (do not change within planning horizon)
\item Historical data provides accurate estimates of travel time distributions
\item Road reliability can be modeled probabilistically
\end{enumerate}

\textbf{Operational Assumptions:}
\begin{enumerate}
\item Service time at each node is constant or negligible
\item Loading/unloading times are included in travel time estimates
\item Vehicles can only depart from warehouses (no en-route rerouting)
\item No vehicle breakdowns or driver unavailability
\item Fuel is available at all warehouses
\end{enumerate}

\textbf{Assumption Justification:}

These assumptions are justified as follows:
\begin{itemize}
\item \textit{Network connectivity:} Somalia road network data shows sufficient connectivity for our study region
\item \textit{Homogeneous vehicles:} Humanitarian organizations typically use standardized vehicle fleets
\item \textit{Known demands:} Health centers provide supply requests in advance
\item \textit{Independent travel times:} Road condition factors (weather, damage) are localized
\item \textit{No en-route rerouting:} Communication infrastructure in Somalia is limited; static plans preferred
\end{itemize}

In Section \ref{sec:sensitivity}, we test the robustness of our solutions to violations of these assumptions.

\subsection{Mathematical Preliminaries}

The problem combines elements from several classical optimization problems:

\textbf{Vehicle Routing Problem (VRP):} Our base structure follows the capacitated VRP (CVRP), a well-known NP-hard problem \cite{Toth2014}.

\textbf{Stochastic Programming:} We employ chance-constrained programming and recourse models to handle uncertainty \cite{Birge2011}.

\textbf{Robust Optimization:} We design solutions that perform well across scenarios, not just on average \cite{Bertsimas2004}.

\textbf{Multi-Objective Optimization:} We consider efficiency-reliability-cost trade-offs using Pareto optimization \cite{Deb2002}.

The problem complexity arises from:
\begin{itemize}
\item Combinatorial explosion: $O(N!)$ possible route orderings
\item Stochastic constraints: Probability calculations require numerical integration or simulation
\item Multi-depot coordination: Inter-warehouse decisions create additional coupling
\end{itemize}

This motivates our algorithmic approach in Section \ref{sec:solution_algorithm}, which combines exact methods for small instances with metaheuristics for realistic-scale problems.
